\paragraph*{Formulae referenced in the problems}

\noindent\textbf{1.1.1}\quad
\begin{equation}
\vec a\times\vec b=
\begin{pmatrix}
a_2b_3-a_3b_2\\
a_3b_1-a_1b_3\\
a_1b_2-a_2b_1
\end{pmatrix}.
\tag{1.1.1.1}\label{eq:cross-product}
\end{equation}
\noindent\textbf{1.1.1}\quad
\begin{equation}
\vec a\times\vec b\perp\vec a,\qquad
\vec a\times\vec b\perp\vec b.
\tag{1.1.1.3}\label{eq:cross-perpendicular}
\end{equation}

\noindent\textbf{1.1.5}\quad
\begin{equation}
\vec a\cdot\vec b=a_1b_1+a_2b_2+a_3b_3,\qquad
\vec a\perp\vec b\iff\vec a\cdot\vec b=0.
\tag{1.1.5.1}\label{eq:dot-product}
\end{equation}

\noindent\textbf{1.1.6}\quad
\begin{equation}
d(\text{parallel lines})=
\frac{\norm{(\vec b-\vec a)\times\vec v}}{\norm{\vec v}}.
\tag{1.1.6.1}\label{eq:parallel-distance}
\end{equation}

\noindent\textbf{1.1.10}\quad
\begin{align}
\text{line parallel to plane}&\iff \vec d\cdot\vec n=0.
\tag{1.1.10.1}\label{eq:line-plane-parallel}\\
\text{plane through }\vec r_0\text{ with normal }\vec n:\quad
\vec n\cdot(\vec r-\vec r_0)&=0.
\tag{1.1.10.2}\label{eq:plane-normal}
\end{align}

\noindent\textbf{1.1.15}\quad
\begin{equation}
\vec r=\vec A+t\vec d
\quad\text{for a line through }\vec A\text{ with direction }\vec d.
\tag{1.1.15.1}\label{eq:line-point-direction}
\end{equation}

\noindent\textbf{1.1.20}\quad
\begin{equation}
y'+P(x)y=Q(x),\qquad
\mu(x)=e^{\int P(x)\,dx},\qquad
\frac{d}{dx}(\mu y)=\mu Q.
\tag{1.1.20.1}\label{eq:if-20}
\end{equation}

\noindent\textbf{1.1.21}\quad
\begin{equation}
(1+(y')^2)^3=r^2(y'')^2
\quad\text{for a circle family of fixed radius }r.
\tag{1.1.21.1}\label{eq:circle-family}
\end{equation}

\noindent\textbf{1.1.22}\quad
\begin{equation}
y=vx\quad\Longrightarrow\quad y'=v+xv'
\quad\text{(homogeneous first-order ODE)}.
\tag{1.1.22.1}\label{eq:homogeneous-substitution}
\end{equation}

\noindent\textbf{1.1.24}\quad
\begin{equation}
\mu(x)=e^{\int P(x)\,dx}.
\tag{1.1.24.1}\label{eq:if-24}
\end{equation}

\noindent\textbf{1.1.25}\quad
\begin{equation}
\frac{d}{dx}\log|u(x)|=\frac{u'(x)}{u(x)},\qquad u(x)\ne0.
\tag{1.1.25.1}\label{eq:log-derivative-25}
\end{equation}

\noindent\textbf{1.1.28}\quad
\begin{equation}
\det(AB)=\det(A)\det(B),\qquad
\det(A^n)=(\det A)^n.
\tag{1.1.28.1}\label{eq:det-power}
\end{equation}

\noindent\textbf{1.1.30}\quad
\begin{equation}
A^{-1}=\frac{\operatorname{adj}(A)}{\det A},\qquad \det A\ne0.
\tag{1.1.30.1}\label{eq:inverse-matrix}
\end{equation}

\noindent\textbf{1.1.36}\quad
\begin{equation}
a^2+b^2+c^2-ab-bc-ca
=\frac12\bigl((a-b)^2+(b-c)^2+(c-a)^2\bigr).
\tag{1.1.36.1}\label{eq:sum-squares}
\end{equation}

\noindent\textbf{1.1.37}\quad
\begin{equation}
[A\mid\vec b]\xrightarrow{\text{elementary row operations}}[I\mid\vec x]
\quad\text{(Gaussian elimination)}.
\tag{1.1.37.1}\label{eq:gaussian-elimination}
\end{equation}

\noindent\textbf{1.1.38}\quad
\begin{align}
\int\frac{du}{a^2+u^2}&=\frac1a\tan^{-1}\left(\frac{u}{a}\right)+C
\tag{1.1.38.1}\label{eq:scaled-arctan-integral-38}\\
\int\frac{du}{1+u^2}&=\tan^{-1}u+C.
\tag{1.1.38.2}\label{eq:arctan-integral-38}
\end{align}

\noindent\textbf{1.1.39}\quad
\begin{equation}
\int\frac{du}{a^2-u^2}
=\frac1{2a}\log\left|\frac{a+u}{a-u}\right|+C.
\tag{1.1.39.1}\label{eq:int-rational-39}
\end{equation}

\noindent\textbf{1.1.40}\quad
\begin{equation}
\frac{d}{dx}\log|u(x)|=\frac{u'(x)}{u(x)},\qquad u(x)\ne0.
\tag{1.1.40.1}\label{eq:log-derivative-40}
\end{equation}

\noindent\textbf{1.1.41}\quad
\begin{equation}
\frac{P(x)}{(x-a)(x^2+b^2)}
=Q(x)+\frac{A}{x-a}+\frac{Bx+C}{x^2+b^2}.
\tag{1.1.41.1}\label{eq:partial-fractions-41}
\end{equation}

\noindent\textbf{1.1.42}\quad
\begin{equation}
\frac{d}{dx}\log|u(x)|=\frac{u'(x)}{u(x)},\qquad u(x)\ne0.
\tag{1.1.42.1}\label{eq:log-derivative-42}
\end{equation}

\noindent\textbf{1.1.43}\quad
\begin{equation}
\frac{P(x)}{(x-a)(x^2+b^2)}
=Q(x)+\frac{A}{x-a}+\frac{Bx+C}{x^2+b^2}.
\tag{1.1.43.1}\label{eq:partial-fractions-43}
\end{equation}

\noindent\textbf{1.1.44}\quad
\begin{equation}
\frac{5\sin x+3\cos x}{\sin x+\cos x}
=4+\frac{\sin x-\cos x}{\sin x+\cos x}.
\tag{1.1.44.1}\label{eq:trig-split-44}
\end{equation}

\noindent\textbf{1.1.46}\quad
\begin{equation}
A=\int_a^b|y_{\mathrm{upper}}-y_{\mathrm{lower}}|\,dx.
\tag{1.1.46.1}\label{eq:area-between-46}
\end{equation}

\noindent\textbf{1.1.47}\quad
\begin{equation}
A=\int_a^b|y_{\mathrm{upper}}-y_{\mathrm{lower}}|\,dx.
\tag{1.1.47.1}\label{eq:area-between-47}
\end{equation}

\noindent\textbf{1.1.48}\quad
\begin{equation}
\lim_{x\to a^-}f(x)=f(a)=\lim_{x\to a^+}f(x),\qquad
f'_{-}(a)=f'_{+}(a)=f'(a).
\tag{1.1.48.1}\label{eq:continuity-differentiability-48}
\end{equation}

\noindent\textbf{1.1.49}\quad
\begin{equation}
\begin{aligned}
\text{reflexive: }&(a,a)\in R,\\
\text{symmetric: }&(a,b)\in R\Rightarrow(b,a)\in R,\\
\text{transitive: }&(a,b),(b,c)\in R\Rightarrow(a,c)\in R.
\end{aligned}
\tag{1.1.49.1}\label{eq:relation-properties-49}
\end{equation}

\noindent\textbf{1.1.50}\quad
\begin{equation}
\text{For componentwise addition on }\mathbb R^2:\quad
e=(0,0),\quad(a,b)^{-1}=(-a,-b).
\tag{1.1.50.1}\label{eq:componentwise-operation-50}
\end{equation}

\noindent\textbf{1.1.51}\quad
\begin{equation}
a*b=(a+b)\bmod n,\quad
0\text{ is the identity},\quad
(-a)\bmod n\text{ is the additive inverse}.
\tag{1.1.51.1}\label{eq:modular-operation-51}
\end{equation}

\noindent\textbf{1.1.52}\quad
\begin{equation}
\lim_{x\to a^-}f(x)=f(a)=\lim_{x\to a^+}f(x).
\tag{1.1.52.1}\label{eq:continuity-limits-52}
\end{equation}

\noindent\textbf{1.1.53}\quad
\begin{align}
X\sim\operatorname{Bin}(n,p):\quad
P(X=k)&=\binom nkp^k(1-p)^{n-k}.
\tag{1.1.53.1}\label{eq:binomial-53}\\
G(z)&=\sum_jP(D=j)z^j.
\tag{1.1.53.2}\label{eq:walk-pgf}
\end{align}

\noindent\textbf{1.1.54}\quad
\begin{align}
P(A\mid E)&=\frac{P(E\mid A)P(A)}{P(E)}.
\tag{1.1.54.1}\label{eq:bayes-54}\\
P(E)&=\sum_iP(E\mid A_i)P(A_i).
\tag{1.1.54.2}\label{eq:total-probability-54}
\end{align}

\noindent\textbf{1.1.55}\quad
\begin{align}
E[X]&=\sum_x xP(X=x),\quad
\operatorname{Var}(X)=E[X^2]-E[X]^2.
\tag{1.1.55.1}\label{eq:mean-variance-55}\\
P(X=k)&=\frac{\binom Kk\binom{N-K}{n-k}}{\binom Nn}.
\tag{1.1.55.2}\label{eq:hypergeometric-55}
\end{align}

\noindent\textbf{1.1.56}\quad
\begin{equation}
P(X=k)=\binom nkp^k(1-p)^{n-k}.
\tag{1.1.56.1}\label{eq:binomial-56}
\end{equation}

\noindent\textbf{1.1.57}\quad
\begin{equation}
P(E)=\sum_iP(E\mid A_i)P(A_i).
\tag{1.1.57.1}\label{eq:total-probability-57}
\end{equation}

\noindent\textbf{1.1.58}\quad
\begin{align}
P(A\mid E)&=\frac{P(A\cap E)}{P(E)}.
\tag{1.1.58.1}\label{eq:conditional-probability-58}\\
P(A\mid E)&=\frac{P(E\mid A)P(A)}{P(E)}.
\tag{1.1.58.2}\label{eq:bayes-58-2}
\end{align}

\noindent\textbf{1.1.59}\quad
\begin{equation}
\text{Evaluate the objective at each feasible corner; an optimum occurs at a feasible vertex.}
\tag{1.1.59.1}\label{eq:lpp-59}
\end{equation}

\noindent\textbf{1.1.60}\quad
\begin{equation}
\text{LPP optimum}=\text{best objective value among feasible vertices}.
\tag{1.1.60.1}\label{eq:lpp-60}
\end{equation}

\noindent\textbf{1.1.64}\quad
\begin{equation}
x\cos\theta+y\sin\theta=r.
\tag{1.1.64.1}\label{eq:circle-tangent-64}
\end{equation}

\noindent\textbf{1.1.65}\quad
\begin{equation}
\text{ellipse point }(a\cos t,b\sin t),\qquad
A_{\triangle}=\frac12 bh.
\tag{1.1.65.1}\label{eq:ellipse-area-65}
\end{equation}
